\documentclass[aspectratio=169,11pt]{beamer}

\usetheme{Madrid}
\usecolortheme{default}
\usefonttheme{professionalfonts}
\setbeamertemplate{navigation symbols}{}
\setbeamertemplate{footline}[frame number]

\usepackage{amsmath, amssymb, booktabs, array}

\title{Beliefs, Expectations, and Job Search}
\subtitle{Behavioral Labor -- Week 3}
\author{Teaching deck draft}
\date{}

\begin{document}

\begin{frame}
  \titlepage
\end{frame}

\begin{frame}{Central question}
\begin{itemize}
    \item How do beliefs, expectations, perception, and limited information shape job search?
    \item Week 2 treated worker behavior as a preference problem.
    \item Week 3 asks when search behavior is instead driven by subjective beliefs about prospects, wages, competition, or fit.
    \item The discipline remains: object, margin, identifying variation, counterfactual, welfare.
\end{itemize}
\end{frame}

\begin{frame}{Bridge from Week 2}
\small
\begin{tabular}{p{0.30\linewidth} p{0.58\linewidth}}
\toprule
Week 2 object & Week 3 shift \\
\midrule
Present bias and commitment & Beliefs about returns to search effort \\
Reference points and effort & Perceived wage offers and reservation behavior \\
Worker-side margins & Search intensity, applications, duration, targeting \\
Welfare caution & Choices may reflect mistaken or sticky expectations \\
\bottomrule
\end{tabular}
\end{frame}

\begin{frame}{Why job search is behaviorally rich}
\begin{itemize}
    \item Search is repeated, high-stakes, and uncertain.
    \item Feedback is noisy: rejection can mean bad luck, poor targeting, weak skills, or a crowded vacancy.
    \item Workers may misperceive job-finding probabilities, wage offers, competition, or their own skill fit.
    \item Firms also lack information about worker skills.
    \item The result is a two-sided behavioral-information environment.
\end{itemize}
\end{frame}

\begin{frame}{Benchmark search problem}
\[
\max_{s_t,\bar{w}_t}\;
p_t(s_t)\mathbb{E}\left[V^E(w)\mid w\ge \bar{w}_t\right]
+\left(1-p_t(s_t)\right)V^U-c(s_t)
\]

\small
\begin{tabular}{p{0.17\linewidth} p{0.70\linewidth}}
\toprule
Object & Interpretation \\
\midrule
$s_t$ & search effort: time, applications, channels, targeting \\
$\bar{w}_t$ & reservation wage or acceptance threshold \\
$p_t(s_t)$ & perceived job-finding probability generated by effort \\
$c(s_t)$ & effort, time, application, and planning costs \\
\bottomrule
\end{tabular}
\end{frame}

\begin{frame}{Subjective beliefs and search effort}
\[
\max_{s_t,\bar{w}_t}\;
p_t^b(s_t)\mathbb{E}^b\left[V^E(w)\mid w\ge \bar{w}_t\right]
+\left(1-p_t^b(s_t)\right)V^U-c(s_t)
\]

\begin{itemize}
    \item Behavioral wedge can enter through subjective job-finding probabilities.
    \item It can also enter through perceived wage offers or perceived competition.
    \item Preferences may be standard while beliefs are biased, sticky, or incomplete.
    \item The empirical task is to name the belief object and the search margin.
\end{itemize}
\end{frame}

\begin{frame}{Unemployment duration and belief updating}
\begin{itemize}
    \item Elicited job-finding beliefs can predict future exits.
    \item Prediction does not imply correct calibration.
    \item Long unemployment spells produce repeated but ambiguous signals.
    \item Weak updating can keep reservation wages high or search strategies unchanged.
    \item Welfare depends on whether optimism is useful persistence or distorted learning.
\end{itemize}
\end{frame}

\begin{frame}{Information provision to job seekers}
\small
\begin{tabular}{p{0.24\linewidth} p{0.28\linewidth} p{0.34\linewidth}}
\toprule
Treatment & Search margin & Interpretation issue \\
\midrule
Brochure or advice & effort, applications, exits & information versus motivation \\
Personalized guidance & targeting and channel choice & strategy quality versus monitoring \\
Prospect information & reservation behavior and duration & belief updating versus selection \\
\bottomrule
\end{tabular}

\vspace{0.6em}
The Altmann-Falk-Jaeger-Zimmermann design is useful because the outcome is search behavior and employment, but the mechanism still needs diagnostic care.
\end{frame}

\begin{frame}{Two-sided limited information}
\begin{itemize}
    \item Workers lack information about wages, fit, prospects, and sometimes their own skills.
    \item Firms lack information about worker skills, reliability, and match quality.
    \item Certification can change worker beliefs and firm screening at the same time.
    \item The Carranza-Garlick-Orkin-Rankin design is therefore not a pure worker-belief treatment.
    \item Ask: who learns what, who observes the signal, and which margin moves?
\end{itemize}
\end{frame}

\begin{frame}{Wage announcements and perceived competition}
\begin{itemize}
    \item Posted wages raise expected payoffs in the benchmark.
    \item Workers may also infer that high-wage vacancies are more competitive or selective.
    \item Applications can respond to the wage as a price and as a signal.
    \item The Belot-Kircher-Muller design separates wage announcement from hidden job quality concerns.
    \item The search margin is vacancy choice and application behavior.
\end{itemize}
\end{frame}

\begin{frame}{Application costs and selection}
\begin{itemize}
    \item Application costs can be monetary, time-based, cognitive, or psychological.
    \item Small costs can sharply change who applies.
    \item The Abebe-Caria-Ortiz-Ospina design links frictions to talent selection.
    \item Interpretation: liquidity, salience, perceived returns, procrastination, and transport costs can all matter.
\end{itemize}
\end{frame}

\begin{frame}{Methods and identification map}
\scriptsize
\begin{tabular}{p{0.25\linewidth} p{0.28\linewidth} p{0.32\linewidth}}
\toprule
Method & Identifies well & Main caution \\
\midrule
Elicited beliefs panel & reported prospects and updating & unobserved heterogeneity \\
Survey expectations & expected wages and durations & measurement error \\
Information experiment & causal search response & information versus motivation \\
Posting experiment & vacancy and application response & payoff versus competition signal \\
Certification experiment & worker and firm learning & two-sided treatment \\
Structural search model & counterfactuals and welfare & model dependence \\
\bottomrule
\end{tabular}
\end{frame}

\begin{frame}{Research lab}
\begin{itemize}
    \item Reproduce: synthetic Mueller-Spinnewijn-Topa-style belief panel.
    \item Diagnose: predictive beliefs, duration dependence, weak updating, and reservation behavior.
    \item Transfer: synthetic certification design inspired by Carranza-Garlick-Orkin-Rankin.
    \item Optional extension: classify wage announcements as payoff signals and perceived-competition signals.
\end{itemize}
\end{frame}

\begin{frame}{Research frontier}
\begin{itemize}
    \item Measuring beliefs is not enough; designs must separate beliefs from constraints and selection.
    \item Information provision can change beliefs, attention, motivation, and compliance at once.
    \item Two-sided information links behavioral job search to hiring and firm screening.
    \item Welfare depends on whether interventions improve choices, reallocate applications, or shift pressure.
\end{itemize}
\end{frame}

\begin{frame}{Bridge to Week 4}
\begin{itemize}
    \item Week 3: beliefs, expectations, perception, and information in job search.
    \item Week 4: attention, salience, complexity, and policy design.
    \item The common object is perceived choice environment.
    \item The next question is how workers respond when the rule exists but is hard to notice, understand, or act on.
\end{itemize}
\end{frame}

\begin{frame}{Takeaways}
\begin{itemize}
    \item Behavioral job search is not one mechanism.
    \item Beliefs can be predictive and still biased or weakly updated.
    \item Wage postings, certifications, and application costs are search-information objects.
    \item Empirical claims must identify the belief object, search margin, and standard alternatives.
\end{itemize}
\end{frame}

\end{document}
